
El presente análisis evalúa, en forma teórica, la posibilidad de inserción de un sistema robótico terrestre autónomo con capacidades de detección y neutralización de amenazas aéreas de baja cota (drones), con foco en su eventual presentación a las Fuerzas Armadas de la República Argentina.

\section{Cliente objetivo primario en Argentina}

El cliente objetivo primario es el Estado Nacional, en particular el sector defensa, ya que el producto se alinea con funciones de vigilancia, protección de objetivos críticos y apoyo a operaciones militares. Dentro de este marco, el \textbf{Fondo Nacional de la Defensa (FONDEF)} constituye un instrumento relevante para proyectos de reequipamiento y modernización.\footnote{Ley 27.565 (FONDEF), Boletín Oficial de la República Argentina, 2020.}

La Ley 27.565 establece que el FONDEF financia la recuperación, modernización e incorporación de material para las Fuerzas Armadas, y además fija criterios que favorecen.\footnote{Ministerio de Defensa de la Nación, Fondo Nacional de la Defensa (FONDEF), sitio oficial.}
\begin{itemize}
    \item Sustitución de importaciones y desarrollo de proveedores locales.
    \item Innovación tecnológica e industrial.
    \item Acciones de investigación y desarrollo en sector público y privado.
\end{itemize}

Estos criterios son consistentes con un desarrollo nacional que combine hardware, software, simulación y capacidades de integración.

\section{Mercados potenciales para el producto}

\subsection{Sector público argentino}
\begin{itemize}
    \item \textbf{Defensa:} presentación de capacidad al Ministerio de Defensa y Fuerzas Armadas, con foco en vigilancia de instalaciones, protección de bases y apoyo a defensa de punto contra drones.
    \item \textbf{Seguridad interior:} fuerzas federales y organismos de seguridad para resguardo de infraestructura crítica y eventos de alta exposición.\footnote{Decreto 21/2025, Boletín Oficial de la República Argentina, operación contractual para sistema antidrones.}
\end{itemize}

En términos de canal formal de contratación, las adquisiciones estatales se publican y gestionan en la plataforma COMPR.AR,\footnote{COMPR.AR, Portal de Compras Públicas de la República Argentina.} con avisos vinculados en Boletín Oficial, incluyendo procesos de organismos militares.\footnote{Boletín Oficial de la República Argentina, ejemplo de contratación de la Armada Argentina, Licitación Pública 0456/2025.}

Como referencia comparada para compras de capacidades complejas de defensa, los procesos de NSPA (OTAN) muestran estructuras documentales con \textit{Statement of Work}, matriz de evaluación técnica y requisitos de seguridad.\footnote{NATO Support and Procurement Agency (NSPA), eProcurement, oportunidad de soporte C-UAS.}

En el plano regional, existen antecedentes de adquisiciones públicas de drones por fuerzas armadas, como la compra de drones de observación del Ejército de Chile.\footnote{ChileCompra, Licitación ID 3385-5-L124, adquisición de drones de observación por el Ejército de Chile.}

\section{Análisis de costos del prototipo}

Para esta etapa se estimó el costo de los componentes principales del prototipo físico. El análisis se limita al prototipo actual y no incluye costos de ingeniería, mano de obra, integración ni escalado productivo.
Para el costo de energía de impresión 3D se considera un consumo típico de impresora Ender 3 en régimen de trabajo continuo.\footnote{Referencia técnica de consumo típico de impresora Ender 3 en operación de impresión continua.}

\begin{table}[H]
\centering
\small
\begin{tabular}{p{6.1cm} c c c}
\toprule
\textbf{Componente} & \textbf{Cant.} & \textbf{Unitario (ARS)} & \textbf{Total (ARS)} \\
\midrule
Telémetro láser de largo alcance con interfaz digital y kit de envío internacional & 1 & 178.150 & 178.150 \\
Escáner LiDAR 2D de 360° para mapeo y navegación (tipo Slamtec A1M8) & 1 & 280.000 & 280.000 \\
Driver doble puente H para motores DC (tipo L298) & 2 & 3.332 & 6.664 \\
Módulo GNSS/GPS de posicionamiento con antena activa (tipo NEO-6M) & 1 & 12.141 & 12.141 \\
Marcadora neumática de baja energía para pruebas de puntería (plataforma airsoft) & 1 & 476.000 & 476.000 \\
Batería sellada de gel 12 V 7 Ah para tracción y electrónica & 2 & 30.000 & 60.000 \\
Conjunto de motores DC con reductora para tracción diferencial (kit x3) & 1 & 184.800 & 184.800 \\
Motor paso a paso NEMA 23 con reductor para eje de torreta & 1 & 138.600 & 138.600 \\
Microcontrolador de control embebido (formato Nano, MCU RP2040) & 1 & 150.000 & 150.000 \\
Cámara USB HD para adquisición de video en tiempo real & 1 & 29.000 & 29.000 \\
Kit de fijación mecánica (varillas roscadas, tuercas y arandelas) & 1 & 28.747 & 28.747 \\
Filamento PLA para fabricación aditiva (bobinas de 1 kg) & 7 & 18.000 & 126.000 \\
Rodamientos radiales sellados para ejes (pack x10, tipo 6202-2RS) & 1 & 24.000 & 24.000 \\
Consumo eléctrico de impresión 3D (40 h de operación, Ender 3) & 1 & 656 & 656 \\
\midrule
\textbf{Total estimado del prototipo} &  &  & \textbf{1.694.758} \\
\bottomrule
\end{tabular}
\caption{Estimación de costos de componentes principales del prototipo.}
\label{tab:costos_prototipo}
\end{table}